% !TEX root = main.tex
\section{Introduction}
% paragraph 1: problem motivation. Same claim as abstract sentence 1.
Observed difficulty moves when training exposure moves, even if the scored skill stays fixed.
A reported score then mixes the task with the learner's preparation.
That mix matters when a low score is later read as a property of the task.
We keep one skill fixed and change the training mixture on purpose.
The skill is exact addition of two zero-padded three-digit integers, scored by exact equality of the sum.
The addition algorithm and the test items stay fixed inside each comparison.
The open question is how much of the measured ranking is exposure.

% paragraph 2: prior assumption. Rankings are stated without the training history under test.
A common assumption is that item structure fixes the ranking before any training history is named.
Conditional description length, for a fixed machine, input, answer, and task, is one such structural quantity \citep{kolmogorov1968three,grunwald2008algorithmic}.
It does not change across our exposure conditions, because those arguments do not change.
Resource laws instead predict average loss from parameter count and data or compute \citep{kaplan2020scaling,hoffmann2022chinchilla}.
They still do not name which items a particular network has practiced.
Curriculum methods order examples by a chosen difficulty signal \citep{bengio2009curriculum}.
Compositional benchmarks show that success depends on how primitives are recombined between training and test \citep{lake2018generalization,keysers2020measuring}.
The shared gap is a ranking stated without the history being tested.

% paragraph 3: contribution. The ranking follows the mixture; the probe helps only when histories differ.
In these experiments the ranking follows the training mixture.
The same items look hard under one mixture and easy under another, at a matched budget.
A starting-loss probe, measured on held-out items before the later training stage, records how much of that history the network absorbed.
The probe improves held-out forecasts when starting checkpoints have different histories.
It adds little when every checkpoint shares one warmup.
The probe is a measurement of a trained network, not an intrinsic label of an item.

% paragraph 4: instantiation. Nine prespecified studies, monotone ridge, online code length, verifiable reward, headroom bound.
We report nine studies that were frozen before their outcomes.
The networks have two hidden rectified layers, six one-hot input digits, and four digit heads.
Early studies vary carry exposure with a shared architecture and budget.
Later studies match answer histograms, and then also match per-column carry rates.
Forecasts use a constrained monotone ridge on log loss, fit on calibration outcomes and scored on held-out widths and mixtures.
One study scores an online description length along the adaptation stream \citep{rissanen1978modeling,voita2020information}.
One study replaces supervised updates with a 0/1 exact-sum reward and a leave-one-out group-relative update \citep{shao2024deepseekmath}.
\texttt{experiments/archive/09\_headroom\_bounded\_rlvr\_utility\_law/} keeps that reward and caps the utility forecast by remaining headroom.
Table~\ref{tab:gates} states every prespecified decision, including the ones that fail.

% paragraph 5: evaluation. Supported results and failed gates, one claim per sentence.
Answer-matched exposure reverses the no-carry versus two-carry ranking at every tested width.
The two endpoint leads are 25.30 [23.14, 27.26] and 24.58 [22.88, 26.19] percentage points.
After heterogeneous pretraining, adding the starting-loss probe cuts joint held-out root mean squared error by 47.20\% [42.88, 53.08].
With per-column carry rates matched, the same comparison cuts loss error by 36.22\% [32.45, 40.51] and online description-length error by 56.93\% [55.01, 58.67].
A verifiable-reward stage cuts post-training loss error by 39.87\% [36.20, 50.17], and the accuracy gain itself reverses with history.
The failures are part of the same record.
The first balanced-exposure study recovers most of a long-carry gap and still misses its retention gate, because no-carry accuracy falls by 6.08 points.
A later remediation pool recovers its gap and still misses an 85\% accuracy line, at 81.09\% [79.51, 82.88].
An additive utility law worsens pooled held-out error by 15.78\%, with interval [$-35.91$, $3.15$].
The headroom bound in that study is supported, with relative root mean squared error reduction 73.72\% [70.43, 77.26].
Probe features inside the bound are not established, at 8.10\% [2.52, 14.39].
Every interval above is a percentile interval over eight seed blocks, and the decision rule uses the interval rather than a separate test (p-val=n/a), except for one auxiliary sign-flip test named in Section~\ref{sec:ranking}.

% paragraph 6: implications and contributions.
A difficulty report for this task should name the learner, the training history, and the budget.
We do not estimate what share of public-benchmark difficulty is data.
We do not claim that these networks implement one shared internal algorithm.
\begin{itemize}
\item Exposure reverses observed addition difficulty on fixed items, including after answer matching and after per-column carry matching.
\item A starting-loss probe improves held-out loss forecasts when pretraining histories differ, and it stays below a 10\% line when they do not.
\item The reversal also appears in online description length and in verifiable-reward improvement, beside prespecified gates that fail.
\item A headroom bound repairs the failed additive utility law on fresh seeds.
\item Probe features inside that bound stay below the prespecified 10\% line.
\end{itemize}
